\chapter{El pipeline en cuatro pasos}
\label{cap:pipeline}

\begin{intuicion}
Antes de bucear en componentes individuales conviene tener el mapa
completo. Graph Hunter divide el trabajo en cuatro etapas
\emph{secuenciales} que el analista atraviesa de una y que
internamente se ejecutan en pipeline: \emph{ingesta}, \emph{hipótesis},
\emph{matching}, y \emph{scoring + paginación}. Este capítulo
recorre las cuatro de punta a punta, mostrando qué objeto entra y
qué objeto sale, con un walkthrough concreto de ``apretar Run'' en
el frontend para que el lector vea cómo dialogan capas.
\end{intuicion}

\section{Vista de pájaro}

\begin{figure}[H]
\centering
\begin{tikzpicture}[
    node distance=0.6cm and 0.7cm,
    every node/.style={font=\footnotesize, align=center},
    stage/.style={draw=accent, thick, rounded corners=2mm,
                  minimum width=2.6cm, minimum height=1.0cm,
                  fill=accentlight},
    arr/.style={->, >=Latex, accent, thick}
]
\node[stage] (ingest) {Ingesta\\\textit{Sysmon, EVTX, CSV}};
\node[stage, right=of ingest] (graph) {Grafo temporal\\\textit{en RAM + spill}};
\node[stage, right=of graph] (parse) {Parser DSL\\\textit{Hypothesis AST}};
\node[stage, right=of parse] (match) {Matcher\\\textit{Yannakakis $\bigm|$ DFS $\bigm|$ LFTJ}};
\node[stage, right=of match] (score) {Scoring +\\paginación};

\draw[arr] (ingest) -- (graph);
\draw[arr] (graph) -- (parse);
\draw[arr] (parse) -- (match);
\draw[arr] (match) -- (score);

\node[below=0.3cm of ingest, font=\footnotesize\itshape, accent]
    {ev/sec $\Rightarrow$ aristas/sec};
\node[below=0.3cm of graph, font=\footnotesize\itshape, accent]
    {28~GB EVTX viable};
\node[below=0.3cm of match, font=\footnotesize\itshape, accent]
    {3-hop @ $|V|$=10K en 10~ms};
\end{tikzpicture}
\caption{Pipeline alto-nivel. Cada etapa tiene su propio capítulo
de profundización; las cifras debajo son targets de referencia, no
SLAs duros.}
\label{fig:pipeline-alto-nivel}
\end{figure}

Cada etapa tiene un \emph{contrato de entrada} y un \emph{contrato
de salida} bien definidos. Eso permite testar cada una en aislamiento
(\emph{property tests}, \emph{benches} \texttt{criterion}) y permite
también que el motor sea reutilizable desde varios \emph{transports}
(Tauri, CLI, MCP, gateway HTTP) ---ninguno toca el grafo
directamente, todos pasan por la misma \texttt{GraphHunterApi}.

\section{Etapa 1: Ingesta}

\textbf{Entrada}: una fuente de \emph{logs} (Sysmon NDJSON, Sentinel
KQL, EVTX binario, CSV genérico, o un \emph{custom parser}).

\textbf{Salida}: aristas e identidades agregadas al grafo temporal.

La ingesta produce, para cada \emph{row} de la fuente, cero o más
\emph{triples}:

\[
\texttt{(source\_id, source\_type, rel\_type, dest\_id, dest\_type, timestamp, metadata)}
\]

\noindent
y los empuja al grafo vía \texttt{add\_entity} (idempotente,
deduplicado por el \emph{interner}) y \texttt{add\_relation}
(append + \emph{streaming mirror}).

Tres decisiones de diseño que aparecen acá y que la Parte~IV explica
en profundidad:

\begin{itemize}
    \item \textbf{Stream-y, no batch.} Los logs se procesan en
    \emph{batches} chicos (default 50.000 rows, $\sim$10 MB) en
    lugar de leer el archivo entero a RAM. Para 28~GB de EVTX, eso
    significa que el \emph{peak RAM} queda en el orden del batch
    size en lugar del tamaño del archivo
    (capítulo~\ref{cap:ingest-anatomy}).
    \item \textbf{Lazy index build.} La ingesta no construye los
    índices auxiliares (NLF, k-hop reachability) ---los deja
    indicados como ``a construir''. La construcción se difiere al
    primer \emph{hunt} (capítulo~\ref{cap:finalize}).
    \item \textbf{Streaming backend canónico.} Cada
    \texttt{add\_relation} hace \emph{single-write} al
    \emph{streaming backend} per-vértice degree-adaptive con
    auto-spill incremental. El backend es el único edge store
    desde stage 2c part 2; sirve simultáneamente al matcher hot
    path, a los index builds (NLF, khop\_reach), y a consumers
    delta-aware (live-tail UI, telemetría) ---capítulo
    \ref{cap:streaming-mirror}.
\end{itemize}

\section{Etapa 2: Hipótesis (parsing de la DSL)}

\textbf{Entrada}: una expresión DSL (string).

\textbf{Salida}: una \texttt{Hypothesis} estructurada (AST).

El analista escribe algo como:

\begin{lstlisting}
User u -[Auth]-> Host h -[Connect]-> IP {k=2} HAVING count(distinct last) >= 5 WITHIN 30m
\end{lstlisting}

El \emph{parser} (en \filepath{platform/dsl}) convierte ese texto en
un valor de tipo \texttt{Hypothesis}:

\begin{lstlisting}[language=Rust]
pub struct Hypothesis {
    pub name: String,
    pub steps: Vec<HypothesisStep>,    // 2 pasos en este ejemplo
    pub k_simplicity: usize,           // 2
    pub aggregations: Vec<AggClause>,  // count(distinct last) >= 5
    pub window_seconds: Option<i64>,   // Some(1800)
    pub shape: HypothesisShape,        // Chain (sin `;`)
}
\end{lstlisting}

\noindent
La \texttt{Hypothesis} es el contrato canónico que viaja entre
parser y matcher. Cualquier \emph{transport} (Tauri, MCP, HTTP) que
quiera ejecutar un hunt construye una \texttt{Hypothesis} ---ya sea
parseando DSL o construyéndola programáticamente--- y la pasa a
\texttt{run\_hunt}.

La Parte~VII (capítulos \ref{cap:dsl-design}--\ref{cap:ast})
desarrolla la DSL en detalle: gramática \emph{chain} legacy,
extensión \emph{graph} con el separador \texttt{;}, parser flow,
disambiguación de identificadores.

\section{Etapa 3: Matching}

\textbf{Entrada}: una \texttt{Hypothesis} más una referencia al
grafo temporal.

\textbf{Salida}: \texttt{Vec<HuntResult>}, donde cada resultado es
un \emph{path} resuelto al alfabeto humano.

Esta es la etapa con más carne algorítmica del libro. La Parte~V
(capítulos \ref{cap:dfs-baseline}--\ref{cap:planner}) desarrolla
los cuatro operadores físicos. En este vistazo alto-nivel, el flujo
es:

\begin{enumerate}
    \item El \texttt{search\_temporal\_pattern} llama al
    \emph{planner}, que inspecciona la \texttt{Hypothesis} y devuelve
    un \texttt{DispatchPlan}.
    \item Según el plan, dispatcha al operador adecuado: DFS legacy,
    LFTJ-K\(_4\), Yannakakis chain, o Yannakakis tree-walk.
    \item El operador consulta los índices auxiliares (NLF,
    k-hop) para podar \emph{starts} antes de empezar la enumeración.
    \item La enumeración produce resultados; cada uno se convierte
    en un \texttt{HuntResult} (\texttt{Vec<Arc<str>>}).
\end{enumerate}

\subsection{Por qué cuatro operadores y no uno}

La razón es la del capítulo~\ref{cap:joins-y-grafos}: el AGM bound
predice que para algunas \emph{shapes} estructurales (como \(K_4\) o
\(\alpha\)-acíclicos), un algoritmo \emph{worst-case-optimal} es
estrictamente mejor que el DFS naive. Tener varios operadores y un
\emph{planner} que elige permite cobrar ese \emph{lift} cuando la
hipótesis lo permite, sin penalizar al DFS para los casos donde no
hay nada mejor que hacer.

\section{Etapa 4: Scoring y paginación}

\textbf{Entrada}: \texttt{Vec<HuntResult>} (puede ser hasta el
\emph{cap}, default 10.000).

\textbf{Salida}: la misma lista, ranqueada por rareza, paginada para
el frontend.

El \emph{anomaly scorer} (en \filepath{core/graph-engine/src/scoring})
asigna un score a cada \emph{path} basado en:

\begin{itemize}
    \item Frecuencia histórica de cada arista del path
    (\emph{rarity}: aristas vistas pocas veces puntúan más alto).
    \item Distancia a la baseline del usuario / host involucrado.
    \item Combinación lineal con pesos configurables.
\end{itemize}

El frontend pagina sobre la lista ordenada (página de 50, scroll
infinito hacia abajo). El analista ve los más sospechosos primero;
puede expandir cada uno para ver la cadena de evidencia.

\section{Walkthrough: ``apreté Run''}

Para cerrar el capítulo, sigamos un \emph{hunt} completo a través
del stack. El analista escribe la hipótesis del Sysmon W0
(capítulo~\ref{cap:threat-hunting}) y aprieta el botón ``Run'' en el
frontend Tauri.

\begin{figure}[H]
\centering
\begin{tikzpicture}[
    node distance=0.4cm,
    every node/.style={font=\scriptsize, align=center},
    box/.style={draw=accent, thick, rounded corners=1.5mm,
                minimum height=0.7cm, minimum width=4cm,
                fill=accentlight}
]
\node[box] (frontend) {1. \textbf{React frontend} (\textit{App.tsx})\\
                       \scriptsize lee la string DSL del input};
\node[box, below=of frontend] (cmd) {2. \textbf{Tauri command} (\textit{cmd\_run\_hunt})\\
                                     \scriptsize serializa la request, IPC al daemon Rust};
\node[box, below=of cmd] (api) {3. \textbf{GraphHunterApi::run\_hunt}\\
                                \scriptsize endpoint canónico de \texttt{platform/api}};
\node[box, below=of api] (parse) {4. \textbf{parse\_dsl}\\
                                  \scriptsize string $\rightarrow$ \texttt{Hypothesis} AST};
\node[box, below=of parse] (search) {5. \textbf{search\_temporal\_pattern}\\
                                     \scriptsize entry-point del matcher};
\node[box, below=of search] (planner) {6. \textbf{planner.plan(hypothesis)}\\
                                       \scriptsize devuelve \texttt{DispatchPlan}};
\node[box, below=of planner] (dispatch) {7. \textbf{ejecuta operador}\\
                                         \scriptsize DFS, LFTJ, o Yannakakis};
\node[box, below=of dispatch] (results) {8. \textbf{score + cap + paginate}\\
                                         \scriptsize \texttt{Vec<HuntResult>} truncado y ranked};
\node[box, below=of results] (back) {9. \textbf{return up the stack}\\
                                     \scriptsize JSON al frontend, render en tabla};

\foreach \i/\j in {frontend/cmd, cmd/api, api/parse, parse/search,
                   search/planner, planner/dispatch, dispatch/results,
                   results/back} {
    \draw[->, >=Latex, accent, thick] (\i) -- (\j);
}
\end{tikzpicture}
\caption{Las nueve paradas de un \emph{hunt} desde el botón ``Run''
hasta los resultados renderizados. Cada paso es un capítulo, sub-capítulo
o sección de algún libro de GraphHunter; el \emph{cross-ref} explícito
está en cada uno cuando aparezca.}
\label{fig:pipeline-walkthrough}
\end{figure}

\noindent
Notar que las cuatro etapas conceptuales (ingesta, hipótesis,
matching, scoring) ocurren en momentos distintos:

\begin{itemize}
    \item \textbf{Ingesta} ocurrió \emph{antes} ---el analista la
    disparó cuando cargó el dataset---. En el walkthrough no
    aparece.
    \item \textbf{Hipótesis} es el paso 4 (\emph{parse\_dsl}).
    \item \textbf{Matching} son los pasos 5--7.
    \item \textbf{Scoring + paginación} es el paso 8.
\end{itemize}

\noindent
La separación temporal importa: la ingesta paga el costo de leer 28
GB \emph{una vez} y construir el grafo y los índices. El analista
después corre decenas o cientos de hunts sobre el mismo grafo, cada
uno en sub-segundo, sin pagar de nuevo el costo de ingesta. Esa es
la economía que hace al motor viable.

\section{Resumen}

\begin{itemize}
    \item Cuatro etapas: \textbf{ingesta} (logs $\to$ grafo),
    \textbf{hipótesis} (DSL $\to$ AST), \textbf{matching} (AST + grafo
    $\to$ paths), \textbf{scoring + paginación} (paths $\to$ ranked
    list).
    \item Cada etapa tiene contrato de entrada y salida bien
    definidos; eso permite testar y reusar.
    \item El \emph{matching} es la etapa con cuatro operadores
    físicos distintos, elegidos por un \emph{planner} según la
    estructura de la hipótesis.
    \item La ingesta y los matchers están separados temporalmente:
    pagás la ingesta una vez por dataset, los hunts son baratos
    sobre el grafo ya construido.
\end{itemize}

\begin{ejercicios}
\begin{enumerate}
    \item \starred{} Para cada una de las cuatro etapas, identificá
    cuál crate de GraphHunter (\texttt{core}, \texttt{platform},
    \texttt{apps/cli}, \texttt{apps/tauri}) la implementa
    principalmente. ¿Cuál es la única etapa que vive en \texttt{core}?

    \item Imaginá que el analista quiere correr 100 hunts distintos
    sobre el mismo grafo de 28 GB. Si cada hunt tarda 100 ms y la
    ingesta tardó 2 minutos, ¿cuál es el costo amortizado por hunt?
    Discutí: ¿cuándo conviene mantener el grafo vivo en memoria
    versus recargarlo desde \emph{snapshot}?

    \item En el walkthrough de la figura~\ref{fig:pipeline-walkthrough},
    ¿qué pasos involucran cruce de \emph{boundary} de proceso (IPC,
    serialización)? ¿Cuáles son sólo llamadas a función en el mismo
    proceso?

    \item \starred{} Si la \texttt{Hypothesis} produjera, en el
    matcher, exactamente 9999 resultados, ¿hay manera de saber si la
    cap se saturó? ¿Qué cambia el motor en su API para que el
    frontend lo distinga? (Pista: hay un campo en \texttt{HuntResult}
    o en la respuesta agregada.)

    \item Pensá un caso donde uno querría correr la ingesta y el
    matcher \emph{en paralelo} ---un \emph{hunt} disparado mientras
    los logs todavía están entrando--- y discutí qué garantías de
    consistencia debería ofrecer el motor.
\end{enumerate}
\end{ejercicios}
